\documentclass{beamer}
\usetheme{Madrid}

% pdflatex
\usepackage[utf8]{inputenc}
\usepackage[T2A]{fontenc}
\usepackage[russian,english]{babel}
\usepackage{cmap}
% \usepackage{lmodern}

\usepackage{graphicx}
\usepackage{xcolor}
\usepackage{tikz}
\usepackage{amsmath,amssymb,amsthm}
\usepackage{mathtools}
\usepackage{listings}
\usepackage{booktabs}
\usepackage{array}
\usetikzlibrary{arrows.meta,positioning,shapes.geometric,fit,backgrounds,calc,decorations.pathreplacing}

% Footline: show current section (left) and frame numbers (right)
\setbeamertemplate{footline}{%
  \leavevmode%
  \hbox{%
    \begin{beamercolorbox}[wd=.78\paperwidth,ht=2.5ex,dp=1ex,left]{author in head/foot}%
      \hspace*{1em}\usebeamerfont{footline}\insertsectionhead%
    \end{beamercolorbox}%
    \begin{beamercolorbox}[wd=.22\paperwidth,ht=2.5ex,dp=1ex,right]{date in head/foot}%
      \usebeamerfont{footline}\insertframenumber{} / \inserttotalframenumber\hspace*{1em}%
    \end{beamercolorbox}%
  }%
}

\newcommand{\parpause}[1]{\only<+->{#1\par}}
\newenvironment{definitionbox}[1][Определение]{%
  \begin{exampleblock}{#1}
}{%
  \end{exampleblock}
}

\AtBeginSection[]
{
  \begin{frame}
    \frametitle{Содержание}
    \tableofcontents[currentsection]
  \end{frame}
}

\title{Семинар 9. Консенсус, RAFT и FLP-теорема}
\subtitle{Распределенные вычисления}
\author{Георгий Семенов}
\institute{Институт прикладных компьютерных наук \\ Университет ИТМО}
\date{весна 2026}

\begin{document}

\frame{\titlepage}

\begin{frame}
\frametitle{Организационные моменты}

\begin{itemize}
  \item Через неделю дедлайн по третьему домашнему заданию
  \item Проверки 2-го и 3-го домашнего задания проставлю чуть позже в таблицу
\end{itemize}

\end{frame}

%--------------------------------------------------
\section{Синхронные и асинхронные распределенные системы}

\begin{frame}
\frametitle{Синхрония в распределенных системах}

\begin{definitionbox}[Синхронные системы]
  \begin{itemize}
    \item Известно, что любое сообщение либо доставляется за некоторое время C, либо полностью исчезает
    \item Тогда можно разбить выполнение алгоритма на фазы длиной O(C): в начале фазы все процессы отправляют сообщения, потом ждут, а в конце знают, что получили все сообщения, которые только могли прийти
  \end{itemize}
\end{definitionbox}

\begin{definitionbox}[Асинхронные системы]
  \begin{itemize}
    \item Время передачи сообщения сверху не ограничено
    \item Но гарантируется, что если нет отказов, то сообщение рано или поздно дойдёт
  \end{itemize}
\end{definitionbox}

\end{frame}

\begin{frame}
\frametitle{Как узлы воспринимают время?}
\begin{columns}[T]
\column{0.5\textwidth}
\textbf{Синхронная модель}
\begin{itemize}
  \item Задержка сообщений $\leq \delta$ (известно)
  \item Таймаут $=$ надёжный детектор отказов
  \item Примеры: real-time системы, авионика самолета
\end{itemize}
\vspace{0.5em}
\textbf{Асинхронная модель}
\begin{itemize}
  \item Задержка сообщений --- произвольная
  \item «Упала нода» $\equiv$ «нода очень медленная»
  \item Примеры: интернет, облачные кластеры
\end{itemize}

\column{0.48\textwidth}
\begin{tikzpicture}[scale=0.82]
  \draw[-{Stealth},thick] (0,0) -- (5.3,0);
  \node[font=\footnotesize,align=center] at (0.1,-0.55) {Полностью\\синхронная};
  \node[font=\footnotesize,align=center] at (5.0,-0.55) {Полностью\\асинхронная};
  \shade[left color=green!45!white, right color=red!40!white]
    (0,0.12) rectangle (5,0.52);
  \draw[thick,blue!70!black] (1.1,0) -- (1.1,0.65)
    node[above,font=\scriptsize,blue!70!black] {LAN};
  \draw[thick,orange!80!black] (3.2,0) -- (3.2,0.65)
    node[above,font=\scriptsize,orange!80!black,align=center] {Интернет};
  \draw[thick,red!70!black] (4.5,0) -- (4.5,0.65)
    node[above,font=\scriptsize,red!70!black] {WAN};
  \draw[decorate,decoration={brace,amplitude=4pt,mirror},thick]
    (0.4,-1.0) -- (4.6,-1.0);
  \node[font=\scriptsize,align=center] at (2.5,-1.55)
    {Частичная синхронность\\{\tiny Dwork, Lynch, Stockmeyer 1988}};
\end{tikzpicture}
\vspace{0.5em}
\begin{alertblock}{}
\small Реальность $\approx$ \textbf{частичная синхронность}: обычно быстро, но иногда сеть «проваливается»
\end{alertblock}
\end{columns}
\end{frame}

\begin{frame}
\frametitle{Что может пойти не так? Модели отказов}
\begin{center}
\begin{tikzpicture}[
  lvl/.style={draw, rounded corners=4pt, minimum width=8.2cm,
              minimum height=0.68cm, align=center, font=\small},
  scale=0.88, transform shape
]
  \node[lvl, fill=green!20]  at (0,3.0)
    {\textbf{Crash-stop} \quad нода молча останавливается, больше не отвечает};
  \node[lvl, fill=yellow!35] at (0,2.1)
    {\textbf{Crash-recovery} \quad упала и восстановилась (с дисковым журналом)};
  \node[lvl, fill=orange!30] at (0,1.2)
    {\textbf{Omission} \quad нода работает, но теряет часть сообщений};
  \node[lvl, fill=red!25]    at (0,0.3)
    {\textbf{Byzantine} \quad произвольное поведение: ложь, подделка, сговор};
  \draw[-{Stealth},thick,gray] (6.0,3.4) -- (6.0,-0.1)
    node[midway,right,font=\footnotesize,gray,align=left] {сложность \\ алгоритма};
  
\end{tikzpicture}
\end{center}
\vspace{-0.2em}
\begin{block}{Сегодня}
Фокус на \textbf{crash-stop}.
Таймаут --- надёжный детектор отказов \emph{только} в синхронной модели.
\end{block}
\end{frame}

%--------------------------------------------------
\section{Понятие консенсуса}

\begin{frame}
\frametitle{Задача консенсуса}

\begin{definitionbox}[Консенсус]
$n$ процессов, каждый \textbf{предлагает} значение $v_i$.
Все должны \textbf{принять решение} --- одно общее значение.
\end{definitionbox}

\vspace{0.4em}
\textbf{Три свойства:}

\begin{description}
  \item[\textcolor{green!60!black}{Validity}]
    Если все предложили одно $v$ --- решение равно $v$.\\
    \textit{\small Защита от «тупого» алгоритма, который всегда возвращает 0.}
  \item[\textcolor{blue!70!black}{Agreement}]
    Все честные процессы принимают \emph{одинаковое} решение.\\
    \textit{\small Суть задачи: никакого раскола.}
  \item[\textcolor{orange!80!black}{Termination}]
    Каждый честный процесс \emph{в итоге} принимает решение.\\
    \textit{\small Алгоритм не зависает навечно.}
\end{description}

\vspace{0.3em}
\begin{exampleblock}{Заметка}
\small Иногда добавляют \textbf{Integrity}: решение обязано быть значением, которое предложил хотя бы один процесс (не «из воздуха»).
\end{exampleblock}
\end{frame}

\begin{frame}
\frametitle{Где консенсус встречается в жизни?}

\begin{columns}[T]
\column{0.5\textwidth}
\begin{block}{Distributed commit (2PC)}
  Все узлы БД должны решить: \textbf{COMMIT} или \textbf{ABORT}.\\
  \small Если координатор --- единственная точка отказа, 2PC не даёт консенсуса при сбое координатора.
\end{block}
\vspace{0.3em}
\begin{block}{State Machine Replication}
  Реплики должны применять команды в \textbf{одном порядке}.\\
  \small Иначе состояния разойдутся.
\end{block}

\column{0.4\textwidth}
\begin{block}{Leader election}
  Кто координатор прямо сейчас?\\
  \small Два «лидера» одновременно --- катастрофа (split-brain).
\end{block}
\vspace{0.3em}
\begin{block}{Блокчейн}
  Какой блок следующий в цепочке?\\
  \small Proof-of-Work, Proof-of-Stake --- всё это варианты консенсуса.
\end{block}
\end{columns}

\end{frame}

\begin{frame}
\frametitle{Синхронный консенсус: \texorpdfstring{$f{+}1$}{f+1} раундов}

\textbf{Алгоритм} (crash-stop, синхронная модель, $f$ возможных отказов):

\begin{enumerate}
  \item На каждом раунде $r = 1,\ldots,f{+}1$: каждый живой процесс рассылает \emph{всем} все значения, которые он узнал к этому моменту.
  \item После $f{+}1$ раундов: взять $\min$ из всех известных значений (или любую детерминированную функцию).
\end{enumerate}

\vspace{0.2em}
\begin{exampleblock}{Почему $f{+}1$ раундов достаточно?}
В худшем случае каждый раунд «забирает» одну ноду ровно перед отправкой.
За $f$ раундов могут упасть $f$ нод --- но их данные \emph{уже успели разойтись}.
На раунде $f{+}1$ все живые узлы имеют одинаковую картину.
\end{exampleblock}

\vspace{0.2em}
\textbf{Стоимость:} $O(n^2)$ сообщений за раунд $\Rightarrow$ $O(f \cdot n^2)$ всего.\\
\small Дорого при больших $f$, но \emph{корректно}.

\end{frame}

%--------------------------------------------------
\section{Алгоритмы консенсуса}

\begin{frame}
\frametitle{Почему асинхронный консенсус сложен?}

\begin{columns}[T]
\column{0.52\textwidth}
\textbf{Дилемма алгоритма:}

\vspace{0.3em}
\begin{itemize}
  \item \textbf{Ждать все ответы?}\\
    Если нода упала --- зависнем навечно.\\
    \textcolor{red}{Нарушение Termination.}
  \vspace{0.4em}
  \item \textbf{Не ждать всех?}\\
    Разные процессы соберут разные наборы голосов и могут принять \emph{разные} решения.\\
    \textcolor{red}{Нарушение Agreement.}
\end{itemize}

% \vspace{0.5em}
% \begin{alertblock}{Ключевой факт}
% Это \textbf{не сложность реализации} --- это фундаментальное ограничение. (FLP, следующий раздел)
% \end{alertblock}

\column{0.44\textwidth}
\begin{tikzpicture}[node distance=0.7cm, font=\small,
  msg/.style={-{Stealth},thick},
  proc/.style={draw,rounded corners=2pt,minimum width=1.5cm,minimum height=0.45cm,align=center}
]
  \node[proc,fill=blue!15]  (c)  {Клиент};
  \node[proc,fill=green!20,below=0.5cm of c,xshift=-1.75cm] (p1) {Нода 1};
  \node[proc,fill=green!20,below=0.5cm of c]               (p2) {Нода 2};
  \node[proc,fill=red!20,  below=0.5cm of c,xshift=1.75cm]  (p3) {Нода 3\\\tiny (упала?)};

  \draw[msg] (c)  -- (p1);
  \draw[msg] (c)  -- (p2);
  \draw[msg,dashed,red] (c) -- (p3);

  \draw[msg] (p1.south) -- ++(0,-0.35) node[below,font=\scriptsize] {OK};
  \draw[msg] (p2.south) -- ++(0,-0.35) node[below,font=\scriptsize] {OK};
  \node[font=\scriptsize,red] at (p3.south |- {(0,-1.85)}) {?};

  \node[font=\scriptsize,align=center,text=red] at (0,-2.5)
    {Ждать ответа от Нода\,3\\или нет?};
\end{tikzpicture}
\end{columns}
\end{frame}

% \begin{frame}
% \frametitle{Карта алгоритмов консенсуса}

% \begin{center}
% \renewcommand{\arraystretch}{1.25}
% \begin{tabular}{@{}lll@{}}
% \toprule
% \textbf{Алгоритм} & \textbf{Год} & \textbf{Ключевая идея} \\
% \midrule
% Viewstamped Replication & 1988 & Исторически первый, лидер + лог \\
% \textbf{Paxos} & 1989/98 & Классика, кворумы, фазы Prepare/Accept \\
% Multi-Paxos & --- & Paxos для последовательности значений \\
% \textbf{Raft} & 2014 & «Paxos, но понятный»; явный лидер \\
% PBFT & 1999 & Byzantine-отказы; $O(n^2)$ сообщений \\
% Ben-Or & 1983 & Рандомизация монеткой; обходит FLP \\
% \bottomrule
% \end{tabular}
% \end{center}

% \vspace{0.3em}
% \begin{block}{На следующем семинаре}
% Paxos и Raft --- подробно: фазы, edge-кейсы, реализации.\\
% Сегодня: \emph{зачем} они нужны и общая идея --- \textbf{лидер + кворум}.
% \end{block}

% \vspace{0.2em}
% \textbf{Общая схема лидер-based алгоритмов:}
% \begin{enumerate}
%   \item Выбрать лидера (снова консенсус!)
%   \item Лидер собирает кворум и предлагает значение
%   \item Если лидер упал --- выбрать нового
% \end{enumerate}
% \end{frame}

\begin{frame}
\frametitle{Кворумы: большинство как инструмент}

\begin{columns}[T]
\column{0.45\textwidth}
\textbf{Почему $\lfloor n/2 \rfloor + 1$?}

\vspace{0.4em}
Любые два кворума \textbf{пересекаются} хотя бы в одном узле.
Этот узел «помнит» предыдущее решение --- не даёт двум лидерам принять разные значения.

\vspace{0.5em}
\begin{exampleblock}{Аналогия}
У вас 5 копий ключа от сейфа. Чтобы открыть --- нужно 3. Любые две группы по 3 человека \emph{обязательно} пересекаются.
\end{exampleblock}

\column{0.5\textwidth}
\begin{tikzpicture}[scale=0.5,
  nd/.style={circle,draw,minimum size=0.65cm,font=\small,thick}]

  % 5 nodes in a circle
  \foreach \i/\lbl in {90/1, 18/2, 306/3, 234/4, 162/5} {
    \node[nd] (\lbl) at ({2.0*cos(\i)},{2.0*sin(\i)}) {\lbl};
  }

  % quorum A: 1,2,3 -- blue ellipse
  \begin{scope}[on background layer]
    \node[draw=blue!60,fill=blue!10,ellipse,
          fit=(1)(2)(3),inner sep=3pt,rotate=30,
          label={[blue!70,font=\scriptsize]above right:кворум A}] {};
  \end{scope}

  % quorum B: 1,4,5 -- orange ellipse
  \begin{scope}[on background layer]
    \node[draw=orange!70,fill=orange!10,ellipse,
          fit=(1)(4)(5),inner sep=3pt,rotate=-40,
          label={[orange!80,font=\scriptsize]below left:кворум B}] {};
  \end{scope}

  % highlight intersection node
  \node[nd,fill=green!40,thick] at (1) {1};
  \node[font=\scriptsize,green!50!black] at (0,2.7) {пересечение};
  \draw[-{Stealth},green!50!black,thick] (0,2.55) -- (0.15,2.15);
\end{tikzpicture}
\end{columns}

\vspace{0.4em}
\small Работает при $f < n/2$ crash-stop отказов.
\end{frame}

%--------------------------------------------------
\section{FLP-теорема}

\begin{frame}
\frametitle{FLP: невозможность консенсуса}

\begin{alertblock}{Теорема (Fischer, Lynch, Paterson, 1985)}
В \textbf{полностью асинхронной} системе, где хотя бы \textbf{один} процесс может
упасть (crash-stop), не существует \emph{детерминированного} алгоритма консенсуса,
который \emph{всегда завершается}.
\end{alertblock}

\vspace{0.5em}
Иными словами: \textbf{Safety} (Agreement) и \textbf{Liveness} (Termination) ---
нельзя гарантировать одновременно в асинхронной модели при наличии отказов.

\end{frame}

% \begin{frame}
% \frametitle{Интуиция доказательства FLP}

% \textbf{Ключевое понятие:} конфигурация системы (состояния всех процессов + очередь сообщений).

% \vspace{0.3em}
% \begin{description}
%   \item[Univalent-0 / Univalent-1] Исход \emph{уже предрешён}: решение будет 0 (или 1), что бы ни случилось дальше.
%   \item[Bivalent] Оба исхода ещё \emph{возможны} --- система «в подвешенном состоянии».
% \end{description}

% \vspace{0.3em}
% \begin{block}{Идея доказательства (два шага)}
% \begin{enumerate}
%   \item \textbf{Начальная конфигурация бивалентна.}\\
%     \small Если все предложили 0 --- решение 0; если все 1 --- решение 1. Где-то посередине есть «граничная» начальная конфигурация, которая бивалентна.
%   \vspace{0.2em}
%   \item \textbf{Adversary-планировщик всегда держит систему бивалентной.}\\
%     \small Из любой бивалентной конфигурации можно сделать следующий шаг так, чтобы остаться бивалентной --- задержав сообщение «нужного» процесса. Алгоритм никогда не «решает».
% \end{enumerate}
% \end{block}

% \end{frame}

\begin{frame}
\frametitle{FLP на практике: приговор или руководство?}

FLP говорит об \emph{асинхронной} модели с \emph{детерминированным} алгоритмом.
Оба ограничения можно ослабить:

\vspace{0.3em}
\begin{enumerate}
  \item \textbf{Частичная синхронность} (GST-модель, Dwork--Lynch--Stockmeyer)\\
    \small После некоторого момента $\mathit{GST}$ сеть ведёт себя синхронно.\\
    Paxos и Raft завершаются \emph{в периоды стабильности}.

  \vspace{0.3em}
  \item \textbf{Рандомизация} (Ben-Or 1983)\\
    \small «Монетка» делает алгоритм не детерминированным.\\
    Консенсус достигается с вероятностью 1 --- но не в худшем случае.

  \vspace{0.3em}
  \item \textbf{Failure Detectors} (Chandra--Toueg 1996)\\
    \small Добавить оракул, который \emph{иногда неточно} детектирует отказы.\\
    Класс $\Omega$ (eventually perfect leader detector) достаточен для консенсуса.
\end{enumerate}

\end{frame}

\section{Алгоритм консенсуса RAFT}

\begin{frame}
\frametitle{Демо}

Презентация – Мартин Клеппманн

\href{https://thesecretlivesofdata.com/raft/}{https://thesecretlivesofdata.com/raft/}

\end{frame}

\section{Дискуссия}


\begin{frame}
\frametitle{Итоги}

\begin{itemize}
  \item \textbf{Модели систем:} синхронная / асинхронная / частично-синхронная;
    иерархия отказов от crash-stop до Byzantine.
  \item \textbf{Консенсус} --- три свойства: Validity, Agreement, Termination.
    В синхронной модели решается за $f{+}1$ раундов; в асинхронной --- нетривиально.
  \item \textbf{Алгоритмы:} кворумы + лидер --- общая идея Paxos, Raft и других.
    Paxos и Raft подробно --- на следующем семинаре.
  \item \textbf{FLP (1985):} в полностью асинхронной системе детерминированный консенсус
    \emph{невозможен}. Обходы: частичная синхронность, рандомизация, failure detectors.
  % \item \textbf{ABD (1995):} линеаризуемый регистр --- \emph{слабее} консенсуса и решаем
  %   в асинхронной модели. Write использует read-phase + write-phase;
  %   Read добавляет write-back для гарантии линеаризуемости.
\end{itemize}

\end{frame}

\end{document}